\begin{resumo}
\noindent Este relatório analisa o sistema de gestão da inovação da Capim Software, healthtech--fintech fundada em 2021 que integra gestão de clínicas odontológicas (SaaS), parcelamento ao paciente (BNPL) e adquirência (POS) em uma única plataforma. A partir de visita técnica, entrevista estruturada com um gestor de produto da empresa e levantamento de dados secundários, o trabalho caracteriza o modelo de negócio, o ambiente competitivo e a estrutura organizacional da Capim por meio da Estrela de Galbraith, da configuração de Mintzberg e da análise SWOT. O núcleo do relatório, no entanto, é o diagnóstico do processo decisório de inovação da empresa à luz dos frameworks da disciplina PRO3584: a tipologia contingencial de processos de inovação de \textcite{salerno2015}, a cadeia de valor da inovação de \textcite{hansenbirkinshaw2007}, a gestão de portfólio de \textcite{goffinmitchell2010} e o modelo de ambidestria Descoberta--Incubação--Aceleração discutido por \textcite{salernogomes2018}. A análise mostra que a Capim opera um processo híbrido -- combinando um fluxo tradicional incremental ("Garagem 2 $\rightarrow$ Garagem 1 $\rightarrow$ Rua") com elementos de pausa ativa e atividades paralelas para iniciativas radicais -- mas aplica critérios de decisão essencialmente uniformes a projetos de incerteza muito distinta, e não possui uma função formal de geração e descoberta de ideias, o que configura o elo mais fraco de sua cadeia de valor da inovação. Como proposta de aprimoramento, o trabalho recomenda a criação de um hub de descoberta de oportunidades, a adoção de um \textit{learning plan} formal nas fases de maior incerteza e a diferenciação explícita da governança entre inovação incremental e radical, com mecanismos de orçamento e métricas dedicados.

\vspace{\onelineskip}
\noindent\textbf{Palavras-chave}: Gestão da inovação. Processo decisório. Healthtech. Stage-gates. Ambidestria organizacional.
\end{resumo}
